%% Figure fig:square, input from the section that discusses it.
\begin{figure}
\centering
\begin{tikzpicture}[font=\small, >=Stealth, node distance=11mm and 34mm]
  \node (p)  {$p$ at the source};
  \node[right=of p]  (tp) {$T(p)$ at the target};
  \node[below=of p]  (sb) {$I_s(p)$};
  \node[below=of tp] (tb) {$I_t(T(p))$};
  \draw[->, dashed] (p)  -- node[above] {$T$ --- translate (untrusted)} (tp);
  \draw[->, thick] (p)  -- node[left]  {$I_s$ (judge)} (sb);
  \draw[->, thick] (tp) -- node[right] {$I_t$ (judge)} (tb);
  \draw[->, dashed] (tb) -- node[below] {$\lam$ --- carry back (untrusted)} (sb);
  \path (p) -- node[midway] {$\equiv_\pi$} (tb);
\end{tikzpicture}
\Description{A commuting square: the source program is translated
(dashed, untrusted) to the target program; both are interpreted
(solid, judges); the target run is carried back (dashed, untrusted)
and must agree with the source run on the kept observables.}
\caption{The square, the judgment on programs. Dashed arrows are
transports, generated and untrusted; solid arrows are the two
interpreters, the judges. The square is closed empirically, per
corpus program, at admission; the observable renaming $\pi$ is the
projection inside it. Where the translation turns one source frame
into many target frames, a stimulus map and a bound map ride along
and are judged inside the same square against the depths both
interpreters report.}
\label{fig:square}
\end{figure}
